\subsection{\texorpdfstring{\(p\)-自由子集和 \(p\)-基}{p-自由子集和 p-基}}

定义1. --- 设 K 是一个域，L 是 K 的高度≤1的 p-根式扩张。L 的元素的一个族 \((x_i)_{i \in I}\) 称为在 K 上 p-自由的（相应地，称为 L 在 K 上的 p-基），如果 \(x_i \notin K\) 对所有 \(i \in I\) 成立，且由典范单射导出的、从 \(\otimes_{i \in I} K(x_i)\) 到 L

的同态 \(u_{i}:\mathbf{K}(x_{i})\to \mathbf{L}\) 是单射（相应地，双射）。

若 \(a\) 是 \(\mathbf{L} - \mathbf{K}\) 的元素，则它是高度为 1 的 \(p\) -根式元，因而它在 \(\mathbf{K}\) 上的极小多项式是 \(\mathbf{X}^p - a^p\) （V，第24页，命题1）；所以 \(\{1, a, \ldots, a^{p-1}\}\) 是 \(\mathbf{K}(a)\) 在 \(\mathbf{K}\) 上的一组基。设 \((x_i)_{i \in I}\) 是 \(\mathbf{L} - \mathbf{K}\) 的元素的一个族， \(\Lambda\) 为 \(\mathbf{N}^{(I)}\) 中由所有这样的有限支撑族 \(\alpha = (\alpha_i)_{i \in I}\) 组成的子集：其中 \(\alpha_i < p\) 对所有 \(i \in I\) 成立；命题9（III，第471页）表明，元素 \(\otimes_{i \in I} x_i^{\alpha_i}\) 当 \(\alpha\) 取遍

在 \(\Lambda\) 时构成向量空间 \(\otimes_{i\in I}\mathbf{K}(x_i)\) 在 \(\mathbf{K}\) 上的一组基。此外，典范

同态 \(\otimes_{i\in I}\mathbf{K}(x_i)\) 到 \(L\) 的像为 \(\mathbf{K}(x_i)_{i\in I}\) （V，第18页，推论1），我们

于是有以下命题：

命题1.---设 K 是一个域，L 是 K 的高度≤1的 p-根式扩张，且 \(\mathbf{x} = (x_{i})_{i \in I}\) 为 L 的元素的一个族。则向量空间 \(\mathrm{K}(x_{i})_{i \in I}\) 在 K 上由乘积 \(x^{\alpha} = \prod_{i \in I} x_{i}^{\alpha_{i}}\) （ \(\alpha\) 取自 \(\Lambda\) ）生成。对于

\((x_{i})_{i\in I}\) ，要使它是 \(p\) -自由的（相应地，一个 \(p\) -基），充要条件是族 \((\mathbf{x}^{\alpha})_{\alpha \in \Lambda}\) 在 \(\mathbf{K}\) 上线性无关（相应地，是 \(\mathbf{L}\) 在 \(\mathbf{K}\) 上的一组基）。

推论。---设 L' 是域 K 的一个扩张，由两个线性不相交的子扩张 K' 和 L 生成。假设 L 在 K 上是高度 \(\leqslant 1\) 的 p-根式扩张；则 L' 是 K' 的高度 \(\leqslant 1\) 的 p-根式扩张。此外，L 中一族元素在 K 上为 p-自由的（相应地，为 L 在 K 上的 p-基），其充要条件是：它在 K' 上为 p-自由的（相应地，为 L' 在 K' 上的 p-基）。

我们有 \(L^{p} \subset K \subset K'\) 且 \(L' = K'(L)\) ，从而 \(L'^{p} = K'^{p}(L^{p}) \subset K'\) 。换言之， \(L'\) 是高度 \(\leqslant 1\) 的、关于 \(K'\) 的 p-根式扩张。推论的其余断言由命题1以及 V 第14页命题5立即得出。

注记。--- 设 K 是一个域，L 是 K 的高度 \(\leqslant1\) 的 p-根式扩张， \((x_{i})_{i\in I}\) 为 L 的元素的一个族。设 A 为多项式代数 \(\mathrm{K}[(\mathrm{X}_{i})_{i\in I}]\) ， a 为由多项式 \(X_{i}^{p}-x_{i}^{p}\) 生成的 A 的理想， \(\varphi:A\to L\) 为满足 \(\varphi(\mathrm{X}_{i})=x_{i}\) （对每个 \(i\in I\) ）的 K-代数同态。族 \((x_{i})_{i\in I}\) 是 p-自由的，当且仅当 \(\varphi\) 的核等于 a。这是因为 \(\mathrm{K}[(\mathrm{X}_{i})]/\mathfrak{a}\) 可与代数 \(\otimes\mathrm{K}[\mathrm{X}_{i}]/(\mathrm{X}_{i}^{p}-x_{i}^{p})\) 等同。

设 L 是域 K 的高度 \(\leqslant1\) 的 p-根式扩张。L 的子集 S 称为 p-自由的（相应地，p-基），如果由 S 到自身的恒等映射所定义的族是 p-自由的（相应地，p-基）。要使 L 的元素的族 \((x_{i})_{i\in I}\) 是 p-自由的（相应地，p-基），充要条件是：映射 \(i\mapsto x_{i}\) 是 I 到 L 的一个 p-自由子集（相应地，p-基）上的双射。由命题1，p-自由子集的每个子集都是 p-自由的；反之，若 S 是 L 的子集且其所有有限子集都是 p-自由的，则 S 是 p-自由的。最后，L 在 K 上的 p-基是使得 \(\mathrm{L}=\mathrm{K}(\mathrm{B})\) 的 p-自由子集 B。

命题2. --- 设 K 是一个域，L 是 K 的高度≤1的 p-根式扩张，S 是 L 的一个子集。要使 S 是 p-自由的，充要条件是：对 S 的每个异于 S 的子集 T，都有 K(T) ≠ K(S)。

首先假设 S 是 p-自由的，并令 \(T \neq S\) 为 S 的一个子集。用 \(\Lambda\) 记由 0 与 p-1 之间的整数组成的有限支撑族 \(\alpha = (\alpha_{s})_{s \in S}\) 的集合；令 \(\Lambda'\) 为 \(\Lambda\) 中由这样的族 \(\alpha = (\alpha_{s})_{s \in S}\) 组成的子集：其中 \(\alpha_{s} = 0\) 对所有 s 属于 S-T。再令 \(u_{\alpha} = \prod_{s \in S} s^{\alpha_{s}}\) ，其中 \(\alpha \in \Lambda\) 。则（V，第98页，

命题1）族 \((u_{\alpha})_{\alpha \in \Lambda}\) 是 \(\mathbf{K}(\mathbf{S})\) 在 \(\mathbf{K}\) 上的一组基，而子族 \((u_{\alpha})_{\alpha \in \Lambda'}\) 是 \(\mathbf{K}(\mathbf{T})\) 在 \(\mathbf{K}\) 上的一组基。由于 \(\Lambda' \neq \Lambda\) ，故 \(\mathbf{K}(\mathbf{T}) \neq \mathbf{K}(\mathbf{S})\) 。

现在假设 S 不是 \(p\) -自由的。则存在整数 \(n \geqslant 1\) 与 S 的元素序列 \(x_1, \ldots, x_n\) ，使得 \((x_1, \ldots, x_{n-1})\) 是 \(p\) -自由的，而 \((x_1, \ldots, x_n)\) 不是。我们有 \([\mathbf{K}(x_1, \ldots, x_{n-1}) : \mathbf{K}] = p^{n-1}\) 且 \([\mathbf{K}(x_1, \ldots, x_n) : \mathbf{K}] < p^n\) 。由于 \([\mathbf{K}(x_1, \ldots, x_n) : \mathbf{K}]\) 是 \([\mathbf{K}(x_1, \ldots, x_{n-1}) : \mathbf{K}]\) 的倍数，因此

\[
[ \mathrm{K} (x _ {1}, \dots , x _ {n}): \mathrm{K} ] = [ \mathrm{K} (x _ {1}, \dots , x _ {n - 1}): \mathrm{K} ].
\]

由此 \(x_{n} \in \mathbf{K}(x_{1}, \ldots, x_{n-1})\) 。这表明 \(\mathbf{K}(\mathbf{S} - \{x_{n}\}) = \mathbf{K}(\mathbf{S})\) 。

命题3. --- 设 K 是一个域，L 是 K 的高度≤1的 p-根式扩张，S、T 是 L 的两个子集。则以下条件等价：

\begin{enumerate}
\def\labelenumi{\alph{enumi})}
\item
  子集 \(\mathbf{S}\) 为 \(p\) -自由的（相对于 \(\mathbf{K}\) ），且 \(\mathbf{T}\) 为 \(p\) -自由的（相对于 \(\mathbf{K}(\mathbf{S})\) ）。
\item
  \(\mathbf{S} \cap \mathbf{T} = \varnothing\) ，且 \(\mathbf{S} \cup \mathbf{T}\) 为 \(p\) -自由的（相对于 \(\mathbf{K}\) ）。
\end{enumerate}

若 T 在 \(\mathbf{K}(\mathbf{S})\) 上 p-自由，则 \(\mathbf{T} \cap \mathbf{K}(\mathbf{S}) = \varnothing\) ，更不用说 \(S \cap T = \varnothing\) 。因此可以假设 S 与 T 不相交。设 \(\Lambda\) 为 \(\mathbf{N}^{(S \cup T)}\) 中由所有这样的族 \(\alpha = (\alpha_x)_{x \in S \cup T}\) 组成的子集：其中 \(\alpha_x < p\) 对所有 \(x \in S \cup T\) 成立；并以类似方式定义 \(\Lambda'\) （ \(\mathbf{N}^{(S)}\) 的子集）与 \(\Lambda''\) （ \(\mathbf{N}^{(T)}\) 的子集）。我们可以自然地把 \(\mathbf{N}^{(S \cup T)}\) 与 \(\mathbf{N}^{(S)} \times \mathbf{N}^{(T)}\) 等同，从而把 \(\Lambda\) 与 \(\Lambda' \times \Lambda''\) 等同。对 \(\alpha \in \Lambda\) ，令 \(u_\alpha = \prod_{x \in S \cup T} x^{\alpha_x}\) ，并类似地定义 \(u_\beta'\) 与 \(u_\gamma''\) ，其中 \(\beta \in \Lambda'\) 与

\(\gamma \in \Lambda''\) 。我们有 \(u_{\alpha} = u_{\beta}' u_{\gamma}''\) ，其中 \(\alpha = (\beta, \gamma)\) 属于 \(\Lambda = \Lambda' \times \Lambda''\) 。此外， \((u_{\beta}')_{\beta \in \Lambda'}\) 生成向量 K-空间 K(S)（V，第94页，命题1）。要使 S \(\cup\) T 在 K 上为 \(p\) -自由的，充要条件是族 \((u_{\beta}' u_{\gamma}'')_{\beta \in \Lambda', \gamma \in \Lambda''}\) 在 K 上线性无关（V，第 98 页，命题 1）；而这等价于（II，第222页）：族 \((u_{\beta}')_{\beta \in \Lambda'}\) 在 K 上线性无关且族 \((u_{\gamma}'')_{\gamma \in \Lambda''}\) 在 K(S) 上线性无关。于是 \(a\) ）与 \(b\) ）的等价性由命题1（V，第98页）得出。

推论。---设 L 是 K 的高度≤1的 p-根式扩张，M 是 L 的一个子扩张。则 L 是 M 上的高度≤1的 p-根式扩张，且 M 是 K 上的高度≤1的 p-根式扩张。此外，若 B 是 M 在 K 上的 p-基，C 是 L 在 M 上的 p-基，则 B ∩ C = ∅，且 B ∪ C 是 L 在 K 上的 p-基。

设 K 为一个域。族 \((x_{i})_{i\in I}\) 称为 K 的（绝对）p-基，如果它是 K 在 \(K^{p}\) 上的 p-基。对每个整数 \(n\geqslant1\) ，我们用 \(\Lambda(n)\) 表示 \(\mathbf{N}^{(I)}\) 中由所有这样的 \(\alpha=(\alpha_{i})_{i\in I}\) 组成的子集：其中 \(\alpha_{i}<p^{n}\) 对所有 \(i\in I\) 成立。

命题4. --- 设 \(\mathbf{x} = (x_i)_{i \in I}\) 是 K 的一组 \(p\) -基。对每个整数 \(n \geqslant 1\) ，族 \((\mathbf{x}^\alpha)_{\alpha \in \Lambda(n)}\) 是 K 在 \(\mathbf{K}^{p^n}\)

上的一组基。当 n = 1 时，断言归结为命题 1（V，第 98 页）。集合 \(\Lambda(n)\) 由 \(\mathbf{N}^{(I)}\) 中形如 \(\alpha = \beta + p^{n-1}\gamma\) （其中 \(\beta \in \Lambda(n-1)\) ， \(\gamma \in \Lambda(1)\) ）的元素组成；这样的分解是唯一的，且我们有 \(\mathbf{x}^{\alpha} = \mathbf{x}^{\beta}(\mathbf{x}^{\gamma})^{p^{n-1}}\) 。此外，族 \((x_{i}^{p^{n-1}})_{i \in I}\) 显然是 \(K^{p^{n-1}}\) 在 \(K^{p^n}\) 上的一个 p-基，因此族 \((\mathbf{x}^{\gamma})^{p^{n-1}}\) 是 \(K^{p^{n-1}}\) 在 \(K^{p^n}\) 上的一组基。于是利用 II 第222页命题25，由归纳法即得结论。

